
\chapter*{Expos\'e II\\Tangent Bundles -- Lie Algebras}
\addcontentsline{toc}{chapter}{Expos\'e II. Tangent Bundles -- Lie Algebras}
\markboth{EXPOS\'E II. TANGENT BUNDLES -- LIE ALGEBRAS}%
  {EXPOS\'E II. TANGENT BUNDLES -- LIE ALGEBRAS}

\begin{center}
  \textit{by M. Demazure}
\end{center}

We propose in this expos\'e to construct, in the theory of schemes, the
analogues of the tangent bundles and Lie algebras of classical theory. It
will, however, be useful not to restrict ourselves to schemes proper, but
also to consider certain functors on the category of schemes that are not
necessarily representable (for example, the functors
\(\underline{\operatorname{Hom}}\), \(\underline{\operatorname{Norm}}\),
and so forth). As announced in the preceding expos\'e (cf. \SGARefLink{sga3:i:subsection:1-1}{I~1.1}), we shall
identify a scheme with its associated functor.

On the other hand, the constructions presented below go beyond the setting
of scheme theory. They are equally valid, for example, in the theory of
analytic spaces with nilpotent elements, apart from a few changes of detail.

Before beginning this construction, we must give a few general definitions
that supplement those of \SGARefLink{sga3:i:subsection:1:7}{I~1.7}.

\SGARefTarget{sga3:ii:section:1}\section*{1. The functors
  \(\underline{\operatorname{Hom}}_{Z/S}(X,Y)\)}
\addcontentsline{toc}{section}{1. The functors Hom Z/S (X,Y)}

We resume the notation of \SGARefLink{sga3:i:subsection:1-1}{I~1.1}. We identify the category \(\mathscr C\)
with a full subcategory of
\[
  \widehat{\mathscr C}
  =\operatorname{Hom}(\mathscr C^{\circ},\mathbf{Set})
\]
(in particular, we suppress the underlinings that allowed us to distinguish
graphically an object of \(\widehat{\mathscr C}\) from an object of
\(\mathscr C\)).

Consider the following situation: four objects of
\(\widehat{\mathscr C}\), denoted by \(S,X,Y,Z\), the first of which is in
fact an object of \(\mathscr C\), with \(X\) and \(Y\) over \(Z\), and
\(Z\) over \(S\).

\SGAThreeNativeDiagram[0.24\textwidth]
  {native_diagrams/exp2/exp2_p002_XY_over_ZS.tex}
  {The objects \(X\) and \(Y\) over \(Z\), and \(Z\) over \(S\).}
  {fig:exp2-xy-over-zs}
  {LNM 151, p.~44}

\SGAInlineRefTarget{sga3:ii:definition:1-1}\subsection*{\SGARefLink{sga3:ii:definition:1-1}{Definition 1.1}.}
We define an object of \(\widehat{\mathscr C}_{/S}\), denoted
\(\underline{\operatorname{Hom}}_{Z/S}(X,Y)\), by
\begin{equation}
  \underline{\operatorname{Hom}}_{Z/S}(X,Y)(S')
  =\operatorname{Hom}_{Z_{S'}}(X_{S'},Y_{S'})
  =\operatorname{Hom}_{Z}(X\times_S S',Y),
  \tag{1}\label{sga3:ii:section:1:equation:1}
\end{equation}
for every object \(S'\) of \(\mathscr C_{/S}\). One sees at once that
\(\underline{\operatorname{Hom}}_{Z/S}(X,Y)\) is simply the subobject of
\(\underline{\operatorname{Hom}}_S(X,Y)\) formed by the morphisms compatible
with \(p_X\) and \(p_Y\); that is, the kernel of the pair of morphisms
\[
  \underline{\operatorname{Hom}}_S(X,Y)
  \rightrightarrows
  \underline{\operatorname{Hom}}_S(X,Z),
\]
the first defined by composition with \(p_Y\), and the second being the
constant morphism whose ``image'' is \(p_X\).

As in \SGARefLink{sga3:i:subsection:1:7}{I~1.7}, for every object \(T\) of \(\widehat{\mathscr C}\) over
\(S\), there is a natural bijection
\begin{equation}
  \operatorname{Hom}_S
    \bigl(T,\underline{\operatorname{Hom}}_{Z/S}(X,Y)\bigr)
  \simeq \operatorname{Hom}_Z(X\times_S T,Y).
  \tag{2}\label{sga3:ii:section:1:equation:2}
\end{equation}
Moreover, by \SGARefLink{sga3:i:proposition:1-7-1}{I~1.7.1}, if \(E,F\) are objects of
\(\widehat{\mathscr C}\) over \(Z\), then
\[
  \operatorname{Hom}_Z\bigl(E,\underline{\operatorname{Hom}}_Z(F,Y)\bigr)
  \simeq \operatorname{Hom}_Z(E\times_ZF,Y)
  \simeq
  \operatorname{Hom}_Z\bigl(F,\underline{\operatorname{Hom}}_Z(E,Y)\bigr).
\]
Applying this with \(E=X\) and \(F=Z\times_ST\), one obtains natural
bijections, for every object \(T\) of
\(\widehat{\mathscr C}_{/S}\),
\begin{equation}
\begin{aligned}
  \operatorname{Hom}_S
    \bigl(T,\underline{\operatorname{Hom}}_{Z/S}(X,Y)\bigr)
  &\simeq \operatorname{Hom}_Z(X\times_ST,Y)\\
  &\simeq
  \begin{cases}
    \operatorname{Hom}_Z\bigl(Z\times_ST,
      \underline{\operatorname{Hom}}_Z(X,Y)\bigr),\\
    \operatorname{Hom}_Z\bigl(X,
      \underline{\operatorname{Hom}}_Z(Z\times_ST,Y)\bigr).
  \end{cases}
\end{aligned}
\tag{3}\label{sga3:ii:section:1:equation:3}
\end{equation}
These bijections are functorial in \(T\), and hence give isomorphisms of
\(S\)-functors.\footnote{Editor's note: points \SGARefLink{sga3:ii:section:2}{(2)}, \SGARefLink{sga3:ii:section:3}{(3)}, and \SGARefLink{sga3:ii:section:4}{(4)} have been
expanded; in particular, \SGARefLink{fig:exp2-hom-adjunction-triangle}{(4)} will be used in \SGARefLink{sga3:ii:proposition:3-11}{3.11}.}

\SGAThreeNativeDiagram[0.68\textwidth]
  {native_diagrams/exp2/exp2_p002_hom_adjunction_triangle.tex}
  {The functorial isomorphisms (4).}
  {fig:exp2-hom-adjunction-triangle}
  {LNM 151, p.~44}

Let us point out two special cases of the definition. If \(Z=S\), then
\[
  \underline{\operatorname{Hom}}_{S/S}(X,Y)
  =\underline{\operatorname{Hom}}_S(X,Y).
\]

On the other hand, when \(X=Z\), we set
\begin{equation}
  \prod_{Z/S}Y
  =\underline{\operatorname{Hom}}_{Z/S}(Z,Y),
  \tag{5}\label{sga3:ii:section:1:equation:5}
\end{equation}
so that, by definition,
\[
  \left(\prod_{Z/S}Y\right)(S')
  =\operatorname{Hom}_Z(Z\times_SS',Y)
  \simeq \Gamma(Y_{S'}/Z_{S'}).
\]
The functor
\(\prod_{Z/S}:\widehat{\mathscr C}_{/Z}\to
\widehat{\mathscr C}_{/S}\) is right adjoint to the base-change functor
from \(S\) to \(Z\): for every \(S\)-functor \(U\),
\[
  \operatorname{Hom}_S\left(U,\prod_{Z/S}Y\right)
  =\operatorname{Hom}_Z(U\times_SZ,Y).
\]
If \(\mathscr C=\mathbf{Sch}\) and \(Z\) is an \(S\)-scheme, the functor
\(\prod_{Z/S}\) is called the ``Weil restriction of scalars.''
\footnote{Editor's note: the preceding two sentences have been added.}
One also has an isomorphism
\begin{equation}
  \underline{\operatorname{Hom}}_{Z/S}(X,Y)
  \simeq
  \underline{\operatorname{Hom}}_{X/S}(X,Y\times_ZX)
  =\prod_{X/S}(Y\times_ZX),
  \tag{6}\label{sga3:ii:section:1:equation:6}
\end{equation}
which gives in particular, when \(Z=S\), an isomorphism
\begin{equation}
  \underline{\operatorname{Hom}}_S(X,Y)
  \simeq \prod_{X/S}Y_X.
  \tag{7}\label{sga3:ii:section:1:equation:7}
\end{equation}

\SGAInlineRefTarget{sga3:ii:remark:1-2}\subsection*{\SGARefLink{sga3:ii:remark:1-2}{Remark 1.2}.}
The functor
\(Y\mapsto\underline{\operatorname{Hom}}_{Z/S}(X,Y)\) commutes with
products in the following sense: there is a functorial isomorphism
\[
  \underline{\operatorname{Hom}}_{Z/S}(X,Y\times_ZY')
  \simeq
  \underline{\operatorname{Hom}}_{Z/S}(X,Y)
  \times_S
  \underline{\operatorname{Hom}}_{Z/S}(X,Y').
  \tag{*}\label{sga3:ii:section:1:equation:star}
\]
It follows that if \(Y\) is a \(Z\)-group, respectively a \(Z\)-ring,
and so on, then \(\underline{\operatorname{Hom}}_{Z/S}(X,Y)\) is an
\(S\)-group, respectively an \(S\)-ring, and so on.

\SGAInlineRefTarget{sga3:ii:remark:1-3}\subsection*{\SGARefLink{sga3:ii:remark:1-3}{Remark 1.3}.}
Let \(\pi:M\to Y\) be a \(Y\)-functor in \(\mathcal O_Y\)-modules
(cf. \SGARefLink{sga3:i:paragraph:4-3-3-1}{I~4.3.3.1}), and set
\(H=\underline{\operatorname{Hom}}_{Z/S}(X,Y)\). Then
\(\underline{\operatorname{Hom}}_{Z/S}(X,M)\) has a natural
\(\mathcal O_H\)-module structure; more precisely, for every \(H'\) over
\(H\), the set
\[
  \operatorname{Hom}_H
    \bigl(H',\underline{\operatorname{Hom}}_{Z/S}(X,M)\bigr)
\]
has a natural \(\mathcal O(X\times_SH')\)-module structure.
\footnote{Editor's note: this remark has been added; it is useful for
Corollary~3.11.1.}

Indeed, let \(m:M\times_YM\to M\) and
\(\lambda:\mathcal O_Y\times_YM\to M\) denote the morphisms defining
the structures of a \(Y\)-abelian group and an \(\mathcal O_Y\)-module.
Let \(H'\) be an \(S\)-scheme over
\(H=\underline{\operatorname{Hom}}_{Z/S}(X,Y)\); equivalently, one has
been given a \(Z\)-morphism \(f:X\times_SH'\to Y\), which makes
\(X\times_SH'\) a \(Y\)-object. Then
\[
  \operatorname{Hom}_H
    \bigl(H',\underline{\operatorname{Hom}}_{Z/S}(X,M)\bigr)
\]
is the set of \(Z\)-morphisms \(\phi:X\times_SH'\to M\) such that
\(\pi\circ\phi=f\), that is, the set of \(Y\)-morphisms
\(X\times_SH'\to M\).

If \(\phi,\psi\) are two such morphisms, define \(\phi+\psi\) to be the
composite \(Y\)-morphism
\[
  X\times_SH'\xrightarrow{\ \phi\times\psi\ }
  M\times_YM\xrightarrow{\ m\ }M.
\]
This makes \(\underline{\operatorname{Hom}}_{Z/S}(X,M)\) an abelian group
over \(H=\underline{\operatorname{Hom}}_{Z/S}(X,Y)\).

Similarly, if \(a\in\mathcal O(X\times_SH')\), that is, if
\(a:X\times_SH'\to\mathcal O_S\) is an \(S\)-morphism, define
\(a\phi\) to be the composite \(\lambda\circ(a\times\phi)\), where
\(a\times\phi\) is the \(Y\)-morphism from \(X\times_SH'\) to
\(\mathcal O_Y\times_YM\simeq\mathcal O_S\times_SM\) with components
\(a\) and \(\phi\). This gives
\(\operatorname{Hom}_H(H',\underline{\operatorname{Hom}}_{Z/S}(X,M))\)
an \(\mathcal O(X\times_SH')\)-module structure, functorial in the
\(H\)-object \(H'\).

\section*{\SGAInlineRefTarget{sga3:ii:section:2}2. The schemes \(I_S(\mathcal M)\)}
\addcontentsline{toc}{section}{2. The schemes I S (M)}

\SGAInlineRefTarget{sga3:ii:definition:2-1}\subsection*{\SGARefLink{sga3:ii:definition:2-1}{Definition 2.1}.}
Let \(S\) be a scheme and \(\mathcal M\) a quasi-coherent
\(\mathcal O_S\)-module. Denote by
\(\mathscr D_{\mathcal O_S}(\mathcal M)\) the quasi-coherent
\(\mathcal O_S\)-algebra \(\mathcal O_S\oplus\mathcal M\), where
\(\mathcal M\) is regarded as an ideal of square zero. Denote by
\[
  I_S(\mathcal M)=\operatorname{Spec}
    \mathscr D_{\mathcal O_S}(\mathcal M)
\]
the corresponding \(S\)-scheme.
In particular, put
\[
  \mathscr D_{\mathcal O_S}
  =\mathscr D_{\mathcal O_S}(\mathcal O_S),
  \qquad I_S=I_S(\mathcal O_S),
\]
and call these, respectively, the algebra of dual numbers over \(S\) and
the scheme of dual numbers over \(S\).

The assignment \(\mathcal M\mapsto I_S(\mathcal M)\) is a contravariant
functor from quasi-coherent \(\mathcal O_S\)-modules to \(S\)-schemes.
In particular, the morphisms \(0\to\mathcal M\) and
\(\mathcal M\to0\) define, respectively, the structural morphism
\[
  \rho:I_S(\mathcal M)\longrightarrow I_S(0)=S
\]
and a section \(\varepsilon_{\mathcal M}\) of it, called the
\emph{zero section}.

\subsection*{\SGAInlineRefTarget{sga3:ii:subsection:2-1-1}2.1.1.}
Since \(\mathcal M\mapsto I_S(\mathcal M)\) is contravariant, every
\(a\in\operatorname{End}_{\mathcal O_S}(\mathcal M)\) defines an
\(S\)-endomorphism \(a^*\) of \(I_S(\mathcal M)\), and
\[
  1^*=\operatorname{id},\qquad (ab)^*=b^*\circ a^*,\qquad
  0^*=\varepsilon_{\mathcal M}\circ\rho,\qquad
  a^*\circ\varepsilon_{\mathcal M}=\varepsilon_{\mathcal M}.
\]
Consequently, \(I_S(\mathcal M)\) carries a right action of the
multiplicative monoid of
\(\operatorname{End}_{\mathcal O_S}(\mathcal M)\), commuting with the
\(S\)-morphisms \(I_S(\mathcal M)\to I_S(\mathcal M')\) arising from
morphisms \(\mathcal M'\to\mathcal M\); in particular, the operators
\(a^*\) preserve the zero section.\footnote{Editor's note: what follows
expands the original; the numbering \SGARefLink{sga3:ii:subsection:2-1-1}{2.1.1}--\SGARefLink{sga3:ii:subsection:2-1-3}{2.1.3} has also been added.}

For every \(a\in\operatorname{End}_{\mathcal O_S}(\mathcal M)\), every
\(f:S'\to S\), and every \(m\in I_S(\mathcal M)(S')\), write
\(m\cdot a=a^*(m)\). Then
\[
  m\cdot1=m,\qquad (m\cdot a)\cdot b=m\cdot(ab),\qquad
  m\cdot0=\varepsilon_{\mathcal M}(\rho(m)),
\]
and if \(m=\varepsilon_{\mathcal M}(f)\), then \(m\cdot a=m\).

\SGAInlineRefTarget{sga3:ii:remark:2-1-2}\subsection*{\SGARefLink{sga3:ii:remark:2-1-2}{Remark 2.1.2}.}
Formation of \(I_S(\mathcal M)\) commutes with base extension: there are
canonical isomorphisms
\[
  I_S(\mathcal M)_{S'}
  \simeq I_{S'}(\mathcal M\otimes_{\mathcal O_S}\mathcal O_{S'}).
\]
For simplicity write \(I_{S'}(\mathcal M)=I_S(\mathcal M)_{S'}\). More
generally, if \(X\) is an \(S\)-functor, not necessarily representable,
write \(I_X(\mathcal M)=I_S(\mathcal M)\times_SX\).

\subsection*{\SGAInlineRefTarget{sga3:ii:subsection:2-1-3}2.1.3.}
By the preceding discussion, the multiplicative monoid of
\(\mathcal O(S')\) acts on the \(S'\)-scheme
\(I_{S'}(\mathcal M)\), functorially in \(\mathcal M\). In other words,
the \(S\)-scheme \(I_S(\mathcal M)\) is an object with an operator-monoid
structure \(\mathcal O_S\), functorial in \(\mathcal M\). Thus there is
a morphism of \(S\)-schemes
\[
  \lambda:I_S(\mathcal M)\times_S\mathcal O_S
  \longrightarrow I_S(\mathcal M),
\]
satisfying the evident conditions. For every \(S\)-functor \(X\), base
change gives a morphism of \(X\)-functors
\[
  \lambda_X:I_X(\mathcal M)\times_S\mathcal O_S
  \longrightarrow I_X(\mathcal M),
\]
making \(I_X(\mathcal M)\) an object with operator monoid
\(\mathcal O(X)\). Every
\(a\in\mathcal O(X)=\operatorname{Hom}_S(X,\mathcal O_S)\) defines the
\(X\)-endomorphism
\[
  a^*=\lambda_X\circ
  \bigl(\operatorname{id}_{I_X(\mathcal M)}\times(a\circ\operatorname{pr}_X)\bigr).
\]
Explicitly, if \(x\in X(S')\) and
\(m\in I_S(\mathcal M)(S')=I_{S'}(\mathcal M)(S')\), then
\(a(x)=a\circ x\in\mathcal O(S')\) and
\[
  (m,x)\cdot a=(m\cdot a(x),x).
\]
This operation is functorial in \(\mathcal M\) and preserves the zero
section \(\varepsilon_{\mathcal M}:X\to I_X(\mathcal M)\); thus
\(a^*\circ\varepsilon_{\mathcal M}=\varepsilon_{\mathcal M}\).
\footnote{Editor's note: this paragraph has been added; it will be useful
in \SGARefLink{sga3:ii:remark:3-4-2}{3.4.2} and \SGARefLink{sga3:ii:lemma:4-6-2}{4.6.2}.}

The operation is also ``functorial in \(X\)'' in the following sense. If
\(\pi:Y\to X\) is a morphism of \(S\)-functors and
\(u:\mathcal O(X)\to\mathcal O(Y)\) the corresponding ring morphism,
so that \(u(a)=a\circ\pi\), then the following diagram is commutative.

\SGAThreeNativeDiagram[0.30\textwidth]
  {native_diagrams/exp2/exp2_p005_operator_functoriality_square.tex}
  {Functoriality of the operator action under \(Y\to X\).}
  {fig:exp2-operator-functoriality}
  {Polo--Gille re-edition, Expos\'e II}

\subsection*{\SGAInlineRefTarget{sga3:ii:subsection:2-2}2.2.}
Let \(\mathcal M\) and \(\mathcal N\) be two quasi-coherent
\(\mathcal O_S\)-modules. The commutative diagram below induces the
corresponding diagram of \(S\)-schemes on the following page.

\SGAThreeNativeDiagram[0.26\textwidth]
  {native_diagrams/exp2/exp2_p005_module_sum_diamond.tex}
  {The direct-sum diagram of \(\mathcal O_S\)-modules.}
  {fig:exp2-module-sum-diamond}
  {Polo--Gille re-edition, Expos\'e II}

\SGAThreeNativeDiagram[0.42\textwidth]
  {native_diagrams/exp2/exp2_p006_IS_sum_scheme_diagram.tex}
  {The induced commutative diagram of \(S\)-schemes.}
  {fig:exp2-is-sum-schemes}
  {LNM 151, p.~47}

\subsection*{Proposition 2.2.}
For every \(S\)-scheme \(X\), the diagram of functors over \(S\) obtained
by applying \(\underline{\operatorname{Hom}}_S(-,X)\) to the preceding
diagram is cartesian.

\SGAThreeNativeDiagram[0.58\textwidth]
  {native_diagrams/exp2/exp2_p006_Hom_cartesian_diagram.tex}
  {The cartesian diagram of \(\underline{\operatorname{Hom}}_S\)-functors.}
  {fig:exp2-hom-cartesian}
  {LNM 151, p.~47}

Indeed, for every \(S'\to S\), one must verify that the diagram of sets
obtained by evaluating the functors on \(S'\) is cartesian. Since formation
of \(I_S(\mathcal P)\) commutes with base extension, it is enough to take
\(S'=S\), and hence to verify that the following diagram of sets is
cartesian.

\SGAThreeNativeDiagram[0.58\textwidth]
  {native_diagrams/exp2/exp2_p006_points_cartesian_diagram.tex}
  {The cartesian diagram after evaluation on \(S\).}
  {fig:exp2-points-cartesian}
  {LNM 151, p.~47}

If \(x\in X(S)\), it follows from SGA~1, III~5.1 that
\(X(\varepsilon_{\mathcal M})^{-1}(x)\) is isomorphic, functorially in
\(\mathcal M\), to
\[
  \operatorname{Hom}_{\mathcal O_S}
    \bigl(x^*(\Omega^1_{X/S}),\mathcal M\bigr),
\]
where \(\Omega^1_{X/S}\) denotes the sheaf of relative differentials of
\(X\) over \(S\). This latter functor in \(\mathcal M\) plainly carries a
direct sum of \(\mathcal O_S\)-modules to the product of the corresponding
sets, which proves the proposition.

\SGAInlineRefTarget{sga3:ii:corollary:2-2-1}\subsection*{\SGARefLink{sga3:ii:corollary:2-2-1}{Corollary 2.2.1}.}
Let \(X\) be an \(S\)-scheme and \(\mathcal M\) a finite free
\(\mathcal O_S\)-module. The functor
\(\underline{\operatorname{Hom}}_S(I_S(\mathcal M),X)\), as a functor
over \(X\), is isomorphic to a finite product over \(X\) of copies of
\(\underline{\operatorname{Hom}}_S(I_S,X)\).

\SGAInlineRefTarget{sga3:ii:nota:2-2-2}\subsection*{Nota 2.2.2.}
It follows from the proof of the proposition that
\(\underline{\operatorname{Hom}}_S(I_S,X)\), as an \(X\)-functor, is
isomorphic to \(\mathbf V(\Omega^1_{X/S})\) (\SGARefLink{sga3:i:definition:4-6-1}{I~4.6.1}), and is therefore
represented by the vector fibration
\[
  V(\Omega^1_{X/S}).
\]
